% Generated draft: requires mathematical and source review.
\documentclass[11pt]{article}
\usepackage[margin=25mm]{geometry}
\usepackage{fontspec}
\setmainfont{Latin Modern Roman}
\usepackage{amsmath,amssymb,mathtools,unicode-math}
\setmathfont{Latin Modern Math}
\usepackage{enumitem,needspace,xurl,hyperref}
\setlength{\parindent}{0pt}
\setlength{\parskip}{5pt}
\setlength{\emergencystretch}{3em}
\newcommand{\N}{\mathbb{N}}
\newcommand{\Z}{\mathbb{Z}}
\newcommand{\Q}{\mathbb{Q}}
\newcommand{\R}{\mathbb{R}}
\newcommand{\C}{\mathbb{C}}
\newcommand{\F}{\mathbb{F}}
\newcommand{\A}{\mathbb{A}}
\newcommand{\PP}{\mathbb{P}}
\newcommand{\E}{\mathbb{E}}
\newcommand{\eps}{\varepsilon}
\newcommand{\dd}{\,\mathrm{d}}
\begin{document}
\section*{UnsolvedMath ID: 30000568}
\textbf{Source problem code:} catalogue entry
\subsection{Definitions and mathematical statement}
Use the standard bounded commuting-tuple convention. For \(T=(T_1,\ldots,T_n)\) on a topological vector space \(X\), define
\[
\operatorname{Orb}(T,x)=\{T_1^{k_1}\cdots T_n^{k_n}x:(k_1,\ldots,k_n)\in\Z_{\geq0}^n\}.
\]
The tuple is hypercyclic if some \(x\) has \(\overline{\operatorname{Orb}(T,x)}=X\). An operator \(B\) is cyclic if some vector has dense linear span of \(\{B^kx:k\geq0\}\); cyclicity is not the same as hypercyclicity. On a complex Banach space the point spectrum of its adjoint is \(\sigma_p(B^*)=\{\lambda:\exists0\ne\ell\in X^*,\ B^*\ell=\lambda\ell\}\).

The source collects these separate questions. Given a hypercyclic tuple, does there exist a cyclic \(B\) satisfying \(BT_i=T_iB\) for every \(i\)? Can a commuting cyclic \(B\) also be chosen with \(\sigma_p(B^*)=\varnothing\)? Whether this second clause is required for every input tuple or asks only for some such example is not explicit.

Does there exist a tuple and vector with a somewhere-dense but nondense orbit, meaning
\[
\operatorname{int}_X\overline{\operatorname{Orb}(T,x)}\ne\varnothing,
\qquad\overline{\operatorname{Orb}(T,x)}\ne X?
\]
Somewhere dense refers to the closure's interior, not the interior of the countable orbit itself.

Does there exist a hypercyclic \(n\)-tuple of matrices acting on \(\C^n\), with the number of generators exactly the dimension? Separately, does there exist a hypercyclic \(n\)-tuple of hyponormal operators on an infinite-dimensional Hilbert space? Hyponormality means
\[
T_i^*T_i-T_iT_i^*\geq0\qquad(1\leq i\leq n).
\]
The standard commuting assumption does not add double commutation with adjoints. These are existence questions for distinct classes of tuples, not a single tuple required to satisfy all properties simultaneously.

The source does not fix the general ambient Banach or Fr\'echet category, scalar field, tuple-length quantifier, or whether commutativity is part of its definition. The conventional orbit above must be matched to the report; a noncommuting word semigroup would be a different setting. In particular the adjoint-empty question needs an appropriate infinite-dimensional setting and cannot automatically be imposed in the finite-dimensional subquestion. The supplied literature triage checked 2026-08-21 reports no authoritative resolution of the extracted collection, but gives no separate current-status assessment for each clause.
\subsection{Short English statement}
The record asks about cyclic operators in tuple commutants, empty adjoint point spectrum, somewhere-dense nondense joint orbits, matrix tuples with length equal to dimension, and hyponormal hypercyclic tuples. Each clause has a distinct ambient setting and existence quantifier.
\subsection{Sources}
\begin{itemize}
\item Local UnsolvedMath registry, permanent ID 30000568, catalogue entry, \emph{Spectral and Finite-Dimensional Problems for Hypercyclic Tuples}.
\item \emph{Mini-Workshop: Hypercyclicity and Linear Chaos} (2007), Oberwolfach Reports, pp. 2227--2276, DOI \url{https://doi.org/10.4171/owr/2006/37}.
\item Supplied literature triage checked 2026-08-21 cites only that report; no clause-specific theorem, tuple-length range, or ambient-space convention is supplied.
\end{itemize}
\end{document}
